% This text appears at the start of the chapter before the first numbered section.

Part II examined dark matter annihilation signatures in specific, resolved astrophysical structures. Both strategies we explored encountered fundamental limitations. The Galactic Center presents the strongest expected signal, but it remains inextricably entangled with complex diffuse foregrounds and unresolved astrophysical sources. Searching for individual dark matter subhalos among unassociated gamma-ray point sources offers a much cleaner environment, yet is severely limited by the finite sensitivity of the instrument and the large look-elsewhere effect. The inherent ambiguity of the Galactic Center Excess and the rarity of detectable subhalos expose the boundaries of source-by-source identification.

These limitations motivate a shift in strategy. We move from asking whether any specific source is a dark matter subhalo to a broader question: can we recover the statistical properties of the entire faint source population from the collective photon-count distribution of the gamma-ray sky? Below the formal detection threshold lies a vast population of sources that cannot be resolved individually. They collectively imprint their abundance and flux distribution onto the measured pixel counts, and characterizing this population requires inference techniques capable of extracting population-level observables --- most notably the source-count distribution, $dN/dS$.

This chapter introduces the framework required to probe the unresolved sky and presents a deep learning approach to measuring the gamma-ray $dN/dS$. Section~\ref{sec:limits_individual} builds the conceptual argument for transitioning to population studies. Section~\ref{sec:source_count} formalizes the source-count distribution as the primary observable (cf.\ Section~\ref{sec:extragalactic_sources} for the astrophysical context) and reviews analytical approaches to measuring it. Section~\ref{sec:sbi_cnn} details the limitations of these analytical formulations and introduces simulation-based inference (cf.\ Section~\ref{sec:3.2}) via convolutional neural networks on spherical maps as a robust alternative. Finally, Section~\ref{sec:transition_dnds} frames the original research paper that forms the core of this chapter, where we validate our methodology by recovering the extragalactic $dN/dS$ down to a factor of 50 below the \textit{Fermi}-LAT threshold.

The results of this chapter feed into two downstream analyses. The recovered source-count distribution becomes the empirical prior for the probabilistic cataloging framework developed in Chapter~7. The population-level perspective established here also provides the methodological foundation for the cosmological dark matter searches pursued through cross-correlations in Chapter~8.
